%% abtex2-modelo-relatorio-tecnico.tex, v-1.9.2 laurocesar
%% Copyright 2012-2014 by abnTeX2 group at http://abntex2.googlecode.com/ 
%%
%% This work may be distributed and/or modified under the
%% conditions of the LaTeX Project Public License, either version 1.3
%% of this license or (at your option) any later version.
%% The latest version of this license is in
%%   http://www.latex-project.org/lppl.txt
%% and version 1.3 or later is part of all distributions of LaTeX
%% version 2005/12/01 or later.
%%
%% This work has the LPPL maintenance status `maintained'.
%% 
%% The Current Maintainer of this work is the abnTeX2 team, led
%% by Lauro César Araujo. Further information are available on 
%% http://abntex2.googlecode.com/
%%
%% This work consists of the files abntex2-modelo-relatorio-tecnico.tex,
%% abntex2-modelo-include-comandos and abntex2-modelo-references.bib
%%

% ------------------------------------------------------------------------
% ------------------------------------------------------------------------
% abnTeX2: Modelo de Relatório Técnico/Acadêmico em conformidade com 
% ABNT NBR 10719:2011 Informação e documentação - Relatório técnico e/ou
% científico - Apresentação
% ------------------------------------------------------------------------ 
% ------------------------------------------------------------------------

% Documento criado baseado no template disponibilizado pelo
% Prof. Dr. Geraldo Braz Junior, membro do corpo docente do PPGCC - UFMA
% a partir do template "Template TCC Ciência da Computação UFMA" 
% disponibilizado no OverLeaf

\documentclass[
% -- opções da classe memoir --
12pt,				% tamanho da fonte
openright,			% capítulos começam em pág ímpar (insere página vazia caso preciso)
oneside,			% para impressão em verso e anverso. Oposto a oneside
a4paper,			% tamanho do papel. 
english,			% idioma adicional para hifenização
french,				% idioma adicional para hifenização
spanish,			% idioma adicional para hifenização
brazil				% o último idioma é o principal do documento
]{abntex2}
	
\usepackage{lmodern}				% Usa a fonte Latin Modern			
\usepackage[T1]{fontenc}			% Selecao de codigos de fonte.
\usepackage[utf8]{inputenc}			% Codificacao do documento (conversão automática dos acentos)
\usepackage{lastpage}				% Usado pela Ficha catalográfica
\usepackage{indentfirst}			% Indenta o primeiro parágrafo de cada seção.
\usepackage{color}					% Controle das cores
\usepackage{graphicx}				% Inclusão de gráficos e imagens
\usepackage{microtype} 				% para melhorias de justificação
\usepackage{mathtools}				% Matemática
\usepackage{multirow}				% Tabelas
\usepackage[table,xcdraw]{xcolor}	% Tabelas (cores)
\usepackage{pdfpages}
\usepackage{lscape}
\usepackage{float}
\usepackage{makeidx}
\usepackage{amsmath}
\usepackage{algorithm}
\usepackage[noend]{algpseudocode}
\usepackage{array}
\usepackage{gensymb}
\usepackage{float}

\hyphenpenalty=10000

\usepackage{lipsum}				% para geração de dummy text

\usepackage[brazilian,hyperpageref]{backref}	 % Paginas com as citações na bibl
\usepackage[alf, abnt-etal-list=0]{abntex2cite}	% Citações padrão ABNT
\renewcommand{\backrefpagesname}{Citado na(s) página(s):~}

\renewcommand{\backref}{}

\renewcommand*{\backrefalt}[4]{
	\ifcase #1 %
		Nenhuma citação no texto.%
	\or
		Citado na página #2.%
	\else
		Citado #1 vezes nas páginas #2.%
	\fi}%
	
\newcommand{\STAB}[1]{\begin{tabular}{@{}c@{}}#1\end{tabular}}
	
\titulo{Relatório sobre Inclusão Digital e Inteligência Artificial na Educação}
\autor{Nome do Autor}
\local{São Luís - MA}
\data{2024}
\orientador{Prof. Ma. Patrícia Azevedo de Oliveira}
%\coorientador{Prof. Dr. Anselmo Cardoso de Paiva} %se houver
\instituicao{
Especialização em Análise de Dados e Inteligência Artificial
\par 
Universidade Federal do Maranhão
}
\tipotrabalho{Relatório}
\preambulo{Relatório apresentado como requisito parcial para obtenção do pontos na disciplina de Fundamentos e Práticas na EAD, à Especialização em Análise de Dados e Inteligência Artificial.}

% alterando o aspecto da cor azul
\definecolor{blue}{RGB}{41,5,195}

% informações do PDF
\makeatletter
\hypersetup{
     	%pagebackref=true,
		pdftitle={\@title}, 
		pdfauthor={\@author},
    	pdfsubject={\imprimirpreambulo},
	    pdfcreator={LaTeX with abnTeX2},
		pdfkeywords={abnt}{latex}{abntex}{abntex2}{trabalho acadêmico}, 
		colorlinks=true,       		% false: boxed links; true: colored links
    	linkcolor=blue,          	% color of internal links
    	citecolor=blue,        		% color of links to bibliography
    	filecolor=magenta,      		% color of file links
		urlcolor=blue,
		bookmarksdepth=4
}
\makeatother

% O tamanho do parágrafo é dado por:
\setlength{\parindent}{1.3cm}

% Controle do espaçamento entre um parágrafo e outro:
\setlength{\parskip}{0.2cm}  % tente também \onelineskip

\makeindex

\begin{document}

\selectlanguage{brazil}

\frenchspacing 

% ---
% Capa
% ---
\renewcommand{\imprimircapa}{%
\begin{capa}%
\center
\includegraphics[scale=0.25]{imagens/logoufma.png} \\ 
\vspace{1cm}
\ABNTEXchapterfont\Large UNIVERSIDADE FEDERAL DO MARANHÃO \\
\ABNTEXsectionfont\Large Especialização em Análise de Dados e Inteligência Artificial \\

\vspace*{5cm}

\begin{center}
{\ABNTEXsectionfont\Large\imprimirautor} \\
%\vfill
\ABNTEXchapterfont\bfseries\LARGE\imprimirtitulo
\end{center}
\vfill
\large\imprimirlocal \\
\large\imprimirdata
\vspace*{1cm}
\end{capa}
}	

\imprimircapa

\imprimirfolhaderosto

% ---
% inserir o sumario
% ---

\pdfbookmark[0]{\contentsname}{toc}
\tableofcontents*
\cleardoublepage
	

\textual

% Remove o número do capítulo da numeração de seções e subseções
\renewcommand{\thesection}{\arabic{section}}
\renewcommand{\thesubsection}{\arabic{section}.\arabic{subsection}}

\section{Introdução}
Coloque o conteúdo da introdução aqui.

\begin{itemize}
    \item Item 1
    \item Item 2
    \item Item 3
\end{itemize}

\section{Desenvolvimento}

\subsection{Benefícios da IA na EAD}
Coloque o conteúdo do desenvolvimento aqui.


\subsection{Mudança de Práticas na EAD}
Coloque o conteúdo do desenvolvimento aqui.

\subsection{Otimização do Ensino pelos Professores com IA}
Coloque o conteúdo do desenvolvimento aqui.


\subsection{Limites, Possibilidades e Desafios da Implementação da IA na Educação}
Coloque o conteúdo do desenvolvimento aqui.

\section{Considerações Finais}
Coloque o conteúdo das considerações finais aqui.

\section{Referências}
André, C. F., Azevedo, A. B. de, \& Andrade, F. (2023). Inclusão digital e inteligência artificial na educação: avanços, desafios e oportunidades para alunos e professores da Educação Básica à Educação Superior.
% Substitua "Ano", "Nome da Revista", "volume(edição)", "páginas" pelos dados reais do artigo utilizado.
% Inclua outras referências bibliográficas conforme necessário.

\phantompart

\postextual

\end{document}